\documentclass[11pt]{article}
% Shared preamble for the LSM crash-course series.
\usepackage[a4paper,margin=0.9in]{geometry}
\usepackage{amsmath,amssymb,amsthm,mathtools}
\usepackage[dvipsnames]{xcolor}
\usepackage{enumitem}
\usepackage{booktabs}
\usepackage{longtable}
\usepackage{array}
\usepackage{tcolorbox}
\usepackage{fancyhdr}
\usepackage[colorlinks=true,linkcolor=Blue]{hyperref}
\tcbuselibrary{skins,breakable}

\setlist{itemsep=3pt,topsep=4pt}
\pagestyle{fancy}\fancyhf{}
\rhead{\small Linear Statistical Models, B.Stat.\ 3rd yr}
\cfoot{\small\thepage}

\newcommand{\E}{\mathbb{E}}
\newcommand{\Var}{\operatorname{Var}}
\newcommand{\Cov}{\operatorname{Cov}}
\newcommand{\Prob}{\mathbb{P}}
\newcommand{\R}{\mathbb{R}}
\newcommand{\tr}{\operatorname{tr}}
\newcommand{\rank}{\operatorname{rank}}
\newcommand{\Cs}{\mathcal{C}}
\newcommand{\N}{\mathcal{N}}
\newcommand{\Ybar}{\overline{Y}}
\newcommand{\Xbar}{\overline{X}}
\newcommand{\xbar}{\bar{x}}
\newcommand{\iid}{\overset{\text{iid}}{\sim}}
\newcommand{\indep}{\overset{\text{indep}}{\sim}}
\newcommand{\dto}{\overset{d}{\longrightarrow}}
\newcommand{\dequal}{\overset{d}{=}}
\newcommand{\bY}{\mathbf{Y}}
\newcommand{\bX}{\mathbf{X}}
\newcommand{\bZ}{\mathbf{Z}}
\newcommand{\bW}{\mathbf{W}}
\newcommand{\bb}{\boldsymbol{\beta}}
\newcommand{\bmu}{\boldsymbol{\mu}}
\newcommand{\bep}{\boldsymbol{\varepsilon}}
\newcommand{\bal}{\boldsymbol{\alpha}}
\newcommand{\one}{\mathbf{1}}
\newcommand{\zero}{\mathbf{0}}
\newcommand{\I}{\mathbb{I}}
\newcommand{\prf}{\smallskip\noindent\textbf{Proof.}\ }
\renewcommand{\qed}{\hfill$\blacksquare$}

\newtcolorbox{idea}[1]{colback=Blue!4,colframe=Blue!55,title={#1},
  fonttitle=\bfseries,breakable}
\newtcolorbox{confusion}[1]{colback=BrickRed!5,colframe=BrickRed!60,
  title={Common confusion: #1},fonttitle=\bfseries,breakable}
\newtcolorbox{recipe}[1]{colback=ForestGreen!5,colframe=ForestGreen!55,title={#1},
  fonttitle=\bfseries,breakable}
\newtcolorbox{exam}[1]{colback=Orange!6,colframe=Orange!70,title={#1},
  fonttitle=\bfseries,breakable}

\lhead{\small Gauss--Markov Framework --- Crash Course}

\begin{document}

\begin{center}
{\LARGE\bfseries Matrix Algebra and the Gauss--Markov Framework}\\[4pt]
{\large A Crash Course from First Principles}\\[3pt]
{\small Linear Statistical Models --- Lectures 5--6 (3 Aug -- 5 Aug)}
\end{center}
\vspace{2pt}\hrule\vspace{10pt}

\noindent The second of four crash courses. Crash Course 1 fitted a line. Here we do the
general model $\bY=\bX\bb+\bep$ where $\bX$ may be \emph{rank deficient} --- which is
exactly the situation in ANOVA. All the linear algebra you need is built up first, from a
starting point of vectors, matrices and subspaces.

\tableofcontents
\vspace{8pt}\hrule\vspace{10pt}

\section{Why this exists}

In simple linear regression, $\bX^T\bX$ was invertible and $\hat\bb=(\bX^T\bX)^{-1}\bX^T\bY$
was a formula you could just write down. But consider the very first ANOVA-type model:
\[
  Y_{ij}=\mu+\alpha_i+\varepsilon_{ij},\qquad i=1,\dots,r .
\]
Add $c$ to $\mu$ and subtract $c$ from every $\alpha_i$: the model is unchanged. The
parameters are \textbf{not identifiable}, $\bX$ does not have full column rank, and
$(\bX^T\bX)^{-1}$ \emph{does not exist}. Everything in Crash Course 1 breaks.

\begin{idea}{The two questions this course answers}
\begin{enumerate}
  \item If $\bb$ itself cannot be estimated, \emph{what can}? \\
        Answer: exactly those $a^T\bb$ with $a\in \mathcal{R}(\bX)$ --- the row space.
  \item Among all unbiased estimators that are linear in $\bY$, which is best? \\
        Answer: $a^T\hat\bb$, for \emph{any} solution $\hat\bb$ of the normal equations.
        That is the Gauss--Markov theorem.
\end{enumerate}
\end{idea}

\section{The linear algebra toolkit}

\subsection{Column space, row space, rank}

For an $n\times p$ matrix $A$:
\[
  \Cs(A) := \{Av : v\in\R^p\}\subseteq\R^n \quad(\text{column space}),\qquad
  \mathcal{N}(A):=\{v: Av=0\} \quad(\text{null space}).
\]
Write $\mathcal{R}(A):=\Cs(A^T)$ for the row space. The rank $\rho(A)$ is
$\dim\Cs(A)=\dim\mathcal{R}(A)$.

Two facts used constantly:
\begin{itemize}
  \item $\Cs(A^T)=\Cs(A^TA)$, hence also $\rho(A)=\rho(A^TA)$.
  \item $Au=b$ has a solution $\iff b\in\Cs(A)$.
\end{itemize}

\subsection{Projections}

Let $U_1,U_2$ be subspaces of $\R^n$.
\[
  U_1+U_2=\{u_1+u_2\}, \qquad
  U_1\oplus U_2 = U_1+U_2 \text{ when the representation } u=u_1+u_2 \text{ is unique}.
\]
Uniqueness holds iff $U_1\cap U_2=\{0\}$. If $u=u_1+u_2$ uniquely, the map $P(u)=u_2$ is
the \textbf{projection onto $U_2$ along $U_1$}, and it is linear, so $P(u)=Pu$ for a matrix $P$.

\begin{idea}{Characterisation of a projection matrix (TFAE)}
\begin{enumerate}[label=(\roman*)]
  \item $P$ is a projection matrix;
  \item $P$ is \textbf{idempotent}, $P^2=P$;
  \item $\mathcal{N}(P)=\Cs(I-P)$;
  \item $\rho(P)+\rho(I-P)=n$;
  \item $\Cs(P)+\Cs(I-P)$ is a direct sum.
\end{enumerate}
\end{idea}

Two consequences you will use without thinking:
\begin{itemize}
  \item Eigenvalues of an idempotent matrix are $0$ and $1$ only (if $Pv=\lambda v$ then
        $\lambda v=Pv=P^2v=\lambda^2 v$, so $\lambda^2=\lambda$).
  \item Therefore $\tr(P)=\rho(P)$ --- the trace \emph{counts} the dimension of the image.
        This is why traces keep turning into degrees of freedom.
\end{itemize}

\subsection{Orthogonal projections}

Projecting onto $S$ \emph{along $S^\perp$} gives the \textbf{orthogonal} projection.

\begin{idea}{Characterisation of an orthogonal projector (TFAE)}
\begin{enumerate}[label=(\roman*)]
  \item $P$ is an orthogonal projector;
  \item $P=P^TP$;
  \item $P^T=P$ \textbf{and} $P^2=P$ (symmetric \emph{and} idempotent).
\end{enumerate}
\end{idea}
\prf (iii)$\Rightarrow$(ii) is immediate. For (i)$\Rightarrow$(iii): $\Cs(I-P)=S^\perp$ and
$\Cs(P)=S$ are orthogonal, so $P^T(I-P)=0$, giving $P^T=P^TP$; the right side is symmetric,
so $P^T=P$, and then $P=P^2$. \qed

\begin{recipe}{The formula for the orthogonal projector onto $\Cs(\bX)$}
\[
  P_{\bX}=\bX(\bX^T\bX)^{-}\bX^T,
\]
where $(\bX^T\bX)^{-}$ is any symmetric g-inverse (\S2.6). When $\bX$ has full column rank
this is the familiar $\bX(\bX^T\bX)^{-1}\bX^T$. \textbf{$P_\bX$ does not depend on which
g-inverse you pick} --- see \S2.7.
\end{recipe}

\subsection{Sums and differences of projections}

Two projectors can be combined, but only under a condition --- and the condition is
exactly orthogonality of the two subspaces.

\begin{idea}{Proposition (asked as a 10-mark question in 2024)}
Let $P,Q$ be symmetric orthogonal projection matrices. Then
\[
  P+Q \text{ is an orthogonal projection matrix}\iff PQ=O,
\]
and in that case $\rank(P)+\rank(Q)=\rank(P+Q)$.
\end{idea}

\prf $P+Q$ is symmetric automatically, so everything hinges on idempotence:
\[
  (P+Q)^2=P^2+PQ+QP+Q^2=P+Q+(PQ+QP),
  \qquad\text{so}\qquad
  (P+Q)^2=P+Q\iff PQ+QP=O. \tag{$\ast$}
\]

($\Leftarrow$) If $PQ=O$ then $QP=Q^TP^T=(PQ)^T=O$ too, so $(\ast)$ holds.

($\Rightarrow$) Suppose $PQ+QP=O$. Multiply on the \emph{left} by $P$ and use $P^2=P$:
\[
  PQ+PQP=O. \tag{i}
\]
Multiply on the \emph{right} by $P$:
\[
  PQP+QP=O. \tag{ii}
\]
Subtracting, $PQ-QP=O$, i.e.\ $PQ=QP$; substituting back into $(\ast)$ gives $2PQ=O$.

\emph{Rank.} $P+Q$ is now idempotent, and for idempotent matrices trace counts rank, so
\[
  \rank(P+Q)=\tr(P+Q)=\tr(P)+\tr(Q)=\rank(P)+\rank(Q). \qquad\blacksquare
\]

\begin{idea}{The geometry, and the companion result for differences}
$PQ=O$ says $\Cs(P)\perp\Cs(Q)$: for $u=Pa$, $v=Qb$, $u^Tv=a^TPQb=0$. Then
$\Cs(P+Q)=\Cs(P)\oplus\Cs(Q)$ and the ranks add because dimensions of a direct sum add ---
an alternative to the trace argument.

The mirror statement is the one you use constantly later: if $\Cs(Q)\subseteq\Cs(P)$ then
$PQ=QP=Q$, so
\[
  (P-Q)^2=P-PQ-QP+Q=P-Q,
\]
i.e.\ \textbf{$P-Q$ is an orthogonal projector with $\rank(P-Q)=\rank(P)-\rank(Q)$.} That is
precisely why $P_\bX-P_\bZ$ is a projector in the nested-model decomposition of Crash
Course 3.
\end{idea}

\subsection{A non-negative-definite squeeze}

\begin{idea}{Proposition (asked as an 8-mark question in 2025)}
Let $A,B$ be symmetric $n\times n$ with $B$ idempotent. If $A\succeq0$ and
$I_n-A-B\succeq0$, then
\[
  AB=BA=O_n .
\]
\end{idea}

\prf \textbf{Step 1.} $B$ is symmetric and idempotent, hence the orthogonal projector onto
$\Cs(B)$; so $By=y$ for every $y\in\Cs(B)$, and therefore $y^TBy=y^Ty$.

\textbf{Step 2 (the squeeze).} Let $y\in\Cs(B)$. Non-negative definiteness of $I-A-B$ gives
\[
  0\le y^T(I-A-B)y=y^Ty-y^TAy-\underbrace{y^TBy}_{=\,y^Ty}=-y^TAy,
\]
so $y^TAy\le0$; but $A\succeq0$ gives $y^TAy\ge0$. Hence $y^TAy=0$ on all of $\Cs(B)$.

\textbf{Step 3.} $A\succeq0$ symmetric has a square root $A=L^TL$ (take $L=A^{1/2}$ from the
spectral decomposition). Then $0=y^TAy=\|Ly\|^2$ forces $Ly=\zero$, hence $Ay=L^TLy=\zero$.

\textbf{Step 4.} Every column of $B$ lies in $\Cs(B)$, so $AB=O$; transposing,
$BA=(AB)^T=O$. \qed

\begin{confusion}{$y^TAy=0$ does \emph{not} by itself give $Ay=\zero$}
For a general symmetric $A$ it fails: take $A=\begin{pmatrix}0&1\\1&0\end{pmatrix}$ and
$y=e_1$; then $y^TAy=0$ but $Ay=e_2\ne\zero$. Step 3 needs $A\succeq0$ --- which is exactly
why that hypothesis is in the statement. In an exam, say where you use it.
\end{confusion}

\subsection{Generalised inverses}

\begin{idea}{Definition}
$G$ is a \textbf{g-inverse} of $A$ (written $A^-$) if
\[
  AGA=A .
\]
\end{idea}

\textbf{Why the definition is exactly right.} If $b\in\Cs(A)$, so $b=Au$ for some $u$, then
$x=Gb$ solves $Ax=b$:
\[
  A(Gb)=AGAu=Au=b .
\]
So a g-inverse is precisely a device for solving consistent linear systems. Nothing more is
demanded of it, which is why it is not unique.

\textbf{Every matrix has one.} Take a full-rank factorisation $A_{m\times n}=P_{m\times r}Q^T_{r\times n}$
with $r=\rho(A)$. Then $P$ has a left inverse and $Q^T$ a right inverse, and one assembles
$G$ with $AGA=A$.

\textbf{For symmetric $A$, build it from the spectrum.} Write $A=P\Lambda P^T$ with
$PP^T=P^TP=I$, $\Lambda=\text{diag}(\lambda_1,\dots,\lambda_n)$. Try $G=P\Delta P^T$ with
$\Delta=\text{diag}(\delta_1,\dots,\delta_n)$. Then
\[
  AGA=A \iff \Lambda\Delta\Lambda=\Lambda \iff \lambda_i^2\delta_i=\lambda_i\ \ \forall i
  \iff \delta_i=\begin{cases}1/\lambda_i & \lambda_i\neq 0,\\ \text{anything} & \lambda_i=0.\end{cases}
\]
The freedom sits exactly on the null space. Choosing the ``anything'' to be $0$ gives the
\textbf{Moore--Penrose inverse}
\[
  A^{+}=P\Delta P^T,\qquad \delta_i=\begin{cases}1/\lambda_i&\lambda_i\ne0\\ 0&\lambda_i=0\end{cases}
  \qquad\text{i.e.}\quad
  \Lambda=\begin{pmatrix}\Lambda_1&0\\0&0\end{pmatrix}\Rightarrow
  \Delta=\begin{pmatrix}\Lambda_1^{-1}&0\\0&0\end{pmatrix},
\]
which yields the solution of \emph{smallest norm}.

\subsection{The g-inverse facts worth memorising}

\begin{center}
\begin{tabular}{@{}p{7.6cm}p{7.0cm}@{}}
\toprule
Statement & Use \\
\midrule
General solution of $Ax=y$ ($y\in\Cs(A)$) is $A^-y+(I-A^-A)z$, $z\in\R^n$ arbitrary.
  & Describes \emph{all} solutions of the normal equations. \\
$G$ is a g-inverse of $A$ $\iff$ $AG$ is idempotent with $\rho(AG)=\rho(A)$ $\iff$ $GA$ likewise.
  & Lets you verify a candidate. \\
$AG$ is the projection onto $\Cs(A)$ along $\mathcal{N}(AG)$.
  & Connects g-inverses to projections. \\
$A^T(AA^T)^-$ and $(A^TA)^-A^T$ are always g-inverses of $A$. & Constructive. \\
If $\mathcal{R}(A)\subseteq\mathcal{R}(B)$ and $\Cs(C)\subseteq\Cs(B)$, then $AB^-C$ is
  \textbf{invariant} to the choice of $B^-$.
  & \textbf{This is why $P_\bX$ and $a^T\hat\bb$ are well defined.} \\
$A^-$ is \emph{reflexive} if $A^-AA^-=A^-$; \emph{minimum-norm} if $(A^-A)^T=A^-A$;
  \emph{least-squares} if $(AA^-)^T=AA^-$. Moore--Penrose satisfies all three.
  & Vocabulary. \\
\bottomrule
\end{tabular}
\end{center}

\subsection{Building a g-inverse from a full-rank corner}

The construction below is the standard way to \emph{exhibit} a g-inverse of a rank-deficient
matrix, and it is a short exam question in its own right.

\begin{idea}{Proposition (asked as a 6-mark question in 2024)}
Let $A$ be $m\times m$, partitioned as
$A=\begin{pmatrix}A_{11}&A_{12}\\A_{21}&A_{22}\end{pmatrix}$ with $A_{11}$ of size
$r\times r$, $r<m$, and suppose
\[
  \rank(A)=\rank(A_{11})=r .
\]
Then $G=\begin{pmatrix}A_{11}^{-1}&0\\0&0\end{pmatrix}$ is a g-inverse of $A$.
\end{idea}

\prf $A_{11}$ is $r\times r$ of rank $r$, so $A_{11}^{-1}$ exists and $G$ is well defined.
Compute
\[
  AG=\begin{pmatrix}A_{11}&A_{12}\\A_{21}&A_{22}\end{pmatrix}
     \begin{pmatrix}A_{11}^{-1}&0\\0&0\end{pmatrix}
    =\begin{pmatrix}I_r&0\\A_{21}A_{11}^{-1}&0\end{pmatrix},
\]
\[
  AGA=\begin{pmatrix}I_r&0\\A_{21}A_{11}^{-1}&0\end{pmatrix}
      \begin{pmatrix}A_{11}&A_{12}\\A_{21}&A_{22}\end{pmatrix}
     =\begin{pmatrix}A_{11}&A_{12}\\A_{21}&A_{21}A_{11}^{-1}A_{12}\end{pmatrix}.
\]
So $AGA=A$ holds if and only if
\[
  A_{22}=A_{21}A_{11}^{-1}A_{12},
  \qquad\text{i.e.\ iff the Schur complement } A_{22}-A_{21}A_{11}^{-1}A_{12} \text{ vanishes.}
\]
\textbf{And that is exactly what the rank hypothesis says.} The top row block
$[A_{11}\ A_{12}]$ has $r$ rows and rank $r$ (because $A_{11}$ does), while $\rank(A)=r$; so
the row space of $A$ is spanned by those first $r$ rows and every row of the bottom block is
a combination of them: there is an $(m-r)\times r$ matrix $M$ with
\[
  [A_{21}\ \ A_{22}]=M[A_{11}\ \ A_{12}]
  \quad\Longrightarrow\quad A_{21}=MA_{11},\quad A_{22}=MA_{12}.
\]
Since $A_{11}$ is invertible, $M=A_{21}A_{11}^{-1}$, and substituting into the second
equation gives $A_{22}=A_{21}A_{11}^{-1}A_{12}$. \qed

\begin{confusion}{where the hypothesis $\rank(A)=r$ is used}
In exactly one place: forcing the Schur complement to vanish. Drop it and $AG$ is still
idempotent, but $AGA$ differs from $A$ in the bottom-right block, so $G$ is \emph{not} a
g-inverse. Marks are awarded for naming that step, not for the two matrix multiplications.
\end{confusion}

\section{The Gauss--Markov model}

\begin{idea}{The setup}
\[
  \underset{n\times1}{\bY}=\underset{n\times p}{\bX}\ \underset{p\times1}{\bb}+\underset{n\times1}{\bep},
  \qquad \bep\sim \N_n(\zero,\sigma^2 I_n),\qquad \bX \text{ non-random}.
\]
Unknowns: $\bb$ and $\sigma^2$. No assumption that $\rho(\bX)=p$.
\end{idea}

\subsection{Least squares and the normal equations}

$\hat\bb = \arg\min_{\bb}\|\bY-\bX\bb\|^2$. Differentiating,
\[
  \bX^T(\bY-\bX\bb)=0 \quad\Longleftrightarrow\quad \bX^T\bX\bb=\bX^T\bY
  \qquad\text{(\textbf{normal equations}).}
\]

\begin{idea}{The normal equations \emph{always} have a solution}
A system $Au=b$ is solvable iff $b\in\Cs(A)$. Here $b=\bX^T\bY\in\Cs(\bX^T)=\Cs(\bX^T\bX)=\Cs(A)$.
So a solution always exists --- but it is unique only when $\bX^T\bX$ is non-singular, i.e.\
$\rho(\bX)=p$.
\end{idea}

If $\bX^T\bX$ is singular and $\hat\bb$ is a solution, so is $\hat\bb+z$ for every
$z\in\mathcal{N}(\bX^T\bX)$. The general solution is $\hat\bb=(\bX^T\bX)^-\bX^T\bY$ for any
g-inverse.

The fitted values, however, are always unique:
\[
  \hat\bY=\bX\hat\bb=P_{\bX}\bY,\qquad
  \hat\sigma^2=\frac{\|\bY-\hat\bY\|^2}{n-\rho(\bX)} .
\]

\begin{confusion}{$n-p$ or $n-\rho(\bX)$?}
The residual degrees of freedom is $n-\rho(\bX)$, the \emph{rank}, not the number of columns.
In full-rank regression they agree. In ANOVA with $r$ groups written as $\mu+\alpha_i$ there
are $r+1$ columns but rank $r$, so the d.f.\ is $n-r$ --- not $n-r-1$. Getting this wrong
corrupts every MSE, every $t$, and every $F$ in the question.
\end{confusion}

\subsection{Identifiability}

\begin{idea}{Definition and theorem}
A family $\{P_\theta\}$ is \textbf{identifiable} if $P_{\theta_1}=P_{\theta_2}\Rightarrow\theta_1=\theta_2$.

\textbf{Theorem.} The linear model $\bY\sim\N_n(\bX\bb,\sigma^2 I_n)$ is identifiable
$\iff$ $\bX^T\bX$ is non-singular.
\end{idea}
\prf $\N_n(\bX\bb_1,\sigma_1^2I)=\N_n(\bX\bb_2,\sigma_2^2I)$ iff $\bX\bb_1=\bX\bb_2$ and
$\sigma_1^2=\sigma_2^2$ (a normal law is determined by its mean and variance). So
identifiability is exactly the statement ``$\bX\bb_1=\bX\bb_2\Rightarrow\bb_1=\bb_2$'',
i.e.\ $\mathcal{N}(\bX)=\{0\}$, i.e.\ $\rho(\bX)=p$, i.e.\ $\bX^T\bX$ invertible. \qed

\section{Estimability --- the central idea}

\begin{idea}{Definition}
\emph{Estimable} means: an unbiased estimator exists. A linear parametric function
$a^T\bb$ is \textbf{linearly estimable} if there is an $\ell\in\R^n$ with
$\E(\ell^T\bY)=a^T\bb$ for \textbf{all} $\bb$.
\end{idea}

\begin{idea}{The estimability criterion}
\[
  \boxed{\ a^T\bb \text{ is estimable} \iff a\in\mathcal{R}(\bX)=\Cs(\bX^T)=\Cs(\bX^T\bX). }
\]
\end{idea}
\prf $\ell^T\bY$ is unbiased for $a^T\bb$
$\iff \ell^T\bX\bb=a^T\bb$ for all $\bb$
$\iff \ell^T\bX=a^T$
$\iff a=\bX^T\ell$
$\iff a\in\Cs(\bX^T)$. \qed

\begin{recipe}{How to check estimability in an exam}
Two equivalent tests. Pick whichever is easier for the matrix in front of you.
\begin{enumerate}
  \item \textbf{Row-space test.} Ask whether $a^T$ is a linear combination of the rows of
        $\bX$. Write $a^T=c^T\bX$ and try to solve for $c$; if you can, it is estimable.
  \item \textbf{Null-space test.} $a\in\Cs(\bX^T)=\mathcal{N}(\bX)^\perp$, so
        \[
          a^T\bb \text{ estimable} \iff a^Tz=0 \text{ for every } z\in\mathcal{N}(\bX).
        \]
        Find a basis of $\mathcal{N}(\bX)$ once, then testing any $a$ is a dot product.
        \emph{This is usually the faster route.}
\end{enumerate}
\end{recipe}

\subsection{The worked example from the notes}

Take $Y_{ij}=\alpha_i+\gamma_j+\varepsilon_{ij}$, $i,j\in\{1,2\}$, with
$\bb=(\alpha_1,\alpha_2,\gamma_1,\gamma_2)^T$ and rows ordered $Y_{11},Y_{12},Y_{21},Y_{22}$:
\[
  \bX=\begin{pmatrix}1&0&1&0\\1&0&0&1\\0&1&1&0\\0&1&0&1\end{pmatrix},\qquad \rho(\bX)=3 .
\]
Rank is $3$, not $4$: the model is not identifiable and $\bb$ cannot be estimated. Writing
$a^T=c^T\bX$ with $c=(c_1,c_2,c_3,c_4)^T$ gives
\[
  a^T=(c_1+c_2,\ c_3+c_4,\ c_1+c_3,\ c_2+c_4),
\]
so the constraint on $a=(a_1,a_2,a_3,a_4)^T$ is
\[
  \boxed{a_1+a_2=a_3+a_4}.
\]
Check the null-space route as well: $z=(1,1,-1,-1)^T$ satisfies $\bX z=0$, and
$a^Tz=a_1+a_2-a_3-a_4=0$ is the same condition.

\begin{center}
\begin{tabular}{@{}llc@{}}
\toprule
Function & $a$ & Estimable? \\
\midrule
$\alpha_1$ & $(1,0,0,0)$ & \textbf{No} ($1+0\ne0+0$) \\
$\gamma_1$ & $(0,0,1,0)$ & \textbf{No} \\
$\alpha_1+\gamma_1$ & $(1,0,1,0)$ & \textbf{Yes} ($1=1$) \\
$\gamma_1-\gamma_2$ & $(0,0,1,-1)$ & \textbf{Yes} ($0=0$) \\
$\alpha_1-\alpha_2$ & $(1,-1,0,0)$ & \textbf{Yes} ($0=0$) \\
\bottomrule
\end{tabular}
\end{center}

\begin{idea}{The moral, which is the moral of ANOVA}
Individual effects are not estimable; \textbf{differences} of effects and \textbf{contrasts}
are. This is exactly why in one-way ANOVA you can never estimate $\alpha_i$ on its own but
you can always estimate $\alpha_i-\alpha_{i'}=\mu_i-\mu_{i'}$, and more generally
$\sum c_i\alpha_i$ whenever $\sum c_i=0$.
\end{idea}

\subsection{Uniqueness of the estimate of an estimable function}

Let $\hat\bb_1,\hat\bb_2$ both solve $\bX^T\bX\bb=\bX^T\bY$. Estimability gives $a=\bX^T\bX c$
for some $c$, so
\[
  a^T\hat\bb_1-a^T\hat\bb_2=c^T\bX^T\bX(\hat\bb_1-\hat\bb_2)=c^T(\bX^T\bY-\bX^T\bY)=0 .
\]
So although $\hat\bb$ is not unique, the \emph{estimate of an estimable function} is. This
is what makes the whole rank-deficient theory workable.

Also note the unbiasedness: with $a=\bX^T\bX u$,
$a^T\hat\bb=u^T\bX^T\bX\hat\bb=u^T\bX^T\bY$, hence
$\E(a^T\hat\bb)=u^T\bX^T\bX\bb=a^T\bb$.

\section{BLUE, zero functions, and the Gauss--Markov theorem}

\begin{idea}{Definitions}
\begin{itemize}
  \item A \textbf{linear unbiased estimator} (LUE) of $a^T\bb$ is any $\ell^T\bY$ with
        $\E(\ell^T\bY)=a^T\bb$ for all $\bb$.
  \item A LUE $\ell^T\bY$ is the \textbf{BLUE} if $\Var(\ell^T\bY)\le\Var(m^T\bY)$ for every
        LUE $m^T\bY$ of the same $a^T\bb$.
  \item $c^T\bY$ is a \textbf{linear zero function} (LZF) if $\E(c^T\bY)=0$ for all $\bb$,
        i.e.\ $c^T\bX\bb=0\ \forall\bb$, i.e.\ $\bX^Tc=0$.
\end{itemize}
\end{idea}

Geometrically: LZFs are the linear functions of $\bY$ carrying \emph{no information about
$\bb$} --- they are the residual directions, $c\in\mathcal{N}(\bX^T)=\Cs(\bX)^\perp$.

\subsection{The characterisation of a BLUE}

\begin{idea}{Proposition}
A LUE $\ell^T\bY$ is a BLUE $\iff$ $\Cov(\ell^T\bY,c^T\bY)=0$ for \emph{every} LZF $c^T\bY$.
\end{idea}

\prf ($\Leftarrow$) Let $m^T\bY$ be any other LUE. Then $(m-\ell)^T\bY$ has expectation
$a^T\bb-a^T\bb=0$, so it is an LZF. Hence
\[
  \Var(m^T\bY)=\Var(\ell^T\bY+(m-\ell)^T\bY)
  =\Var(\ell^T\bY)+\Var((m-\ell)^T\bY)+2\underbrace{\Cov(\ell^T\bY,(m-\ell)^T\bY)}_{=0}
  \ \ge\ \Var(\ell^T\bY).
\]

($\Rightarrow$) Let $\ell^T\bY$ be BLUE and $c^T\bY$ an LZF. For every $\lambda\in\R$,
$\ell^T\bY+\lambda c^T\bY$ is still a LUE, so
\[
  \Var(\ell^T\bY)\le \Var(\ell^T\bY)+\lambda^2\Var(c^T\bY)+2\lambda\Cov(\ell^T\bY,c^T\bY),
\]
i.e.\ $0\le \lambda^2\Var(c^T\bY)+2\lambda\Cov(\ell^T\bY,c^T\bY)$ for \emph{all} $\lambda$.
If $\Cov\neq0$, choose $\lambda$ small with the opposite sign to the covariance: the linear
term dominates the quadratic one and the right side is negative --- contradiction. Hence
$\Cov(\ell^T\bY,c^T\bY)=0$. \qed

\begin{idea}{The picture behind this proof}
Decompose any LUE as (the BLUE) $+$ (an LZF). The BLUE lies in $\Cs(\bX)$, the LZFs in
$\Cs(\bX)^\perp$. Being uncorrelated with all LZFs means ``having no wasted component in the
noise directions''. \textbf{BLUE = orthogonal projection, again.}
\end{idea}

\subsection{The Gauss--Markov theorem}

\begin{idea}{Gauss--Markov Theorem}
In the model $\bY=\bX\bb+\bep$ with $\E\bep=\zero$, $\Var(\bep)=\sigma^2I_n$: \\
if $a^T\bb$ is estimable, then $a^T\hat\bb$ is its \textbf{unique BLUE}, where $\hat\bb$ is
\emph{any} solution of the normal equations $\bX^T\bX\bb=\bX^T\bY$.
\end{idea}

\prf Estimability gives $a^T=\ell^T\bX$ for some $\ell$, so $a^T\hat\bb=\ell^T\bX\hat\bb$.
Any solution has the form $\hat\bb=(\bX^T\bX)^-\bX^T\bY+z$ with $z\in\mathcal{N}(\bX^T\bX)$;
but $\mathcal{N}(\bX^T\bX)=\mathcal{N}(\bX)$, so $\bX z=0$ and
\[
  \bX\hat\bb=\bX(\bX^T\bX)^-\bX^T\bY = P_\bX\bY .
\]
Therefore $a^T\hat\bb=\ell^TP_\bX\bY$, which is linear in $\bY$, and
\[
  \E(a^T\hat\bb)=\ell^TP_\bX\bX\bb=\ell^T\bX\bb=a^T\bb,
\]
so it is a LUE. For optimality, write $a=\bX^T\bX u$ (possible by estimability), so
$a^T\hat\bb=u^T\bX^T\bX\hat\bb=u^T\bX^T\bY$. Let $h^T\bY$ be any LZF, so $\bX^Th=0$. Then
\[
  \Cov(a^T\hat\bb,\ h^T\bY)=\Cov(u^T\bX^T\bY,h^T\bY)
  =u^T\bX^T\Var(\bY)h=\sigma^2u^T\bX^Th=0 .
\]
By the previous proposition, $a^T\hat\bb$ is the BLUE. Uniqueness follows from \S4.3. \qed

\begin{exam}{This is Homework 2, Question 1}
HW2 Q1 walks exactly this route: (a) $\hat\bmu=P_\bX\bY$; (b) $\theta=\bX^T\bmu$ is
estimable with BLUE $\hat\theta=\bX^T\bY$; (c) if $a^T\theta$ is estimable, its BLUE is
$w^T\hat\theta$. For (b) note $\theta=\bX^T\bX\bb$, so $a=\bX^T\bX$-times-something for each
coordinate --- estimability is automatic --- and the BLUE is
$\bX^T\hat\bmu=\bX^TP_\bX\bY=\bX^T\bY$ since $\bX^TP_\bX=(P_\bX\bX)^T=\bX^T$.
\end{exam}

\section{Generalised least squares: correlated errors}

Suppose $\bY=\bX\bb+\bep$ with $\bep\sim\N(\zero,\Sigma)$, $\Sigma$ positive definite. Now
LSE and MLE part company.

\textbf{LSE} still solves $\bX^T\bX\bb=\bX^T\bY$ --- it does not know about $\Sigma$.

\textbf{MLE} maximises $-\tfrac12(\bY-\bX\bb)^T\Sigma^{-1}(\bY-\bX\bb)-\tfrac12\log|\Sigma|$, so
\[
  \hat\bb^{ML}=\arg\min_{\bb} (\bY-\bX\bb)^T\Sigma^{-1}(\bY-\bX\bb)
  \quad\Longrightarrow\quad
  \bX^T\Sigma^{-1}\bX\bb=\bX^T\Sigma^{-1}\bY .
\]
With $\bX$ of full column rank,
\[
  \hat\bb^{LS}=(\bX^T\bX)^{-1}\bX^T\bY,\qquad
  \hat\bb^{ML}=(\bX^T\Sigma^{-1}\bX)^{-1}\bX^T\Sigma^{-1}\bY .
\]

\begin{idea}{Both are unbiased, but only one is best}
\[
  \Var(a^T\hat\bb^{LS})=a^T(\bX^T\bX)^{-1}(\bX^T\Sigma\bX)(\bX^T\bX)^{-1}a,
  \qquad
  \Var(a^T\hat\bb^{ML})=a^T(\bX^T\Sigma^{-1}\bX)^{-1}a,
\]
and the second is always $\le$ the first. The right version of Gauss--Markov here is: apply
the ordinary theorem to the transformed model
\[
  \Sigma^{-1/2}\bY=\Sigma^{-1/2}\bX\bb+\Sigma^{-1/2}\bep,
\]
whose error has variance $I$. So the BLUE of an estimable $a^T\bb$ is $a^T\hat\bb^{ML}$,
the \textbf{generalised least squares} estimator. Ordinary least squares is BLUE only when
$\Sigma=\sigma^2 I$ (or, more precisely, when $\Cs(\Sigma\bX)\subseteq\Cs(\bX)$).
\end{idea}

\section{Practice, with answers}

\begin{enumerate}[label=\textbf{P\arabic*.}]
\item Show $\Cs(\bX^T)=\Cs(\bX^T\bX)$ and hence $\mathcal{N}(\bX)=\mathcal{N}(\bX^T\bX)$.

\item In the one-way model $Y_{ij}=\mu+\alpha_i+\varepsilon_{ij}$ with $r=3$ groups, find a
      basis for $\mathcal{N}(\bX)$ and use it to decide which of
      $\mu$, $\alpha_1$, $\alpha_1-\alpha_2$, $\mu+\alpha_1$, $\alpha_1+\alpha_2-2\alpha_3$
      are estimable.

\item Prove that $P_\bX=\bX(\bX^T\bX)^-\bX^T$ does not depend on the choice of g-inverse.

\item Give an example of an LZF in the simple linear regression model, and verify directly
      that it is uncorrelated with $\hat\beta_1$.

\item $\bY\sim\N_3\big((\alpha_1,\ \alpha_1+\alpha_2,\ \alpha_1+\alpha_3)^T,\ \sigma^2I_3\big)$.
      Is $\alpha_2-\alpha_3$ estimable? If so, find its BLUE and its variance.
\end{enumerate}

\subsection*{Answers}

\noindent\textbf{P1.} $\Cs(\bX^T\bX)\subseteq\Cs(\bX^T)$ is clear. Conversely if
$\bX^T\bX v=0$ then $v^T\bX^T\bX v=\|\bX v\|^2=0$, so $\bX v=0$; hence
$\mathcal{N}(\bX^T\bX)\subseteq\mathcal{N}(\bX)$, and the reverse inclusion is trivial, so the
two null spaces coincide. Equal null spaces of $p$-column matrices force equal ranks, so
$\rho(\bX^T\bX)=\rho(\bX)=\rho(\bX^T)$, and the inclusion of column spaces above must be
equality.

\smallskip\noindent\textbf{P2.} With $\bb=(\mu,\alpha_1,\alpha_2,\alpha_3)^T$, every row of
$\bX$ has the form $(1,e_i^T)$, so $\bX z=0$ forces $z=c(1,-1,-1,-1)^T$. Then $a^T\bb$ is
estimable iff $a_0-a_1-a_2-a_3=0$ where $a=(a_0,a_1,a_2,a_3)$:
$\mu$ (i.e.\ $a=(1,0,0,0)$) gives $1\ne0$, \textbf{not estimable};
$\alpha_1$ gives $-1\ne0$, \textbf{not estimable};
$\alpha_1-\alpha_2$ gives $0$, \textbf{estimable};
$\mu+\alpha_1$ gives $1-1=0$, \textbf{estimable} (it is $\mu_1$);
$\alpha_1+\alpha_2-2\alpha_3$ gives $-(1+1-2)=0$, \textbf{estimable} (a contrast).

\smallskip\noindent\textbf{P3.} Apply the invariance rule: $AB^-C$ is invariant to $B^-$ when
$\mathcal{R}(A)\subseteq\mathcal{R}(B)$ and $\Cs(C)\subseteq\Cs(B)$. Take $A=\bX$,
$B=\bX^T\bX$, $C=\bX^T$. Then $\mathcal{R}(\bX)=\Cs(\bX^T)=\Cs(\bX^T\bX)=\mathcal{R}(B)$
and $\Cs(\bX^T)=\Cs(\bX^T\bX)=\Cs(B)$, so both conditions hold with equality.

\smallskip\noindent\textbf{P4.} Any $c$ with $\bX^Tc=0$, i.e.\ $\sum c_i=0$ and
$\sum c_ix_i=0$. For instance with $n=3$ and $x=(1,2,3)$, take $c=(1,-2,1)$. Then
$\Cov(c^T\bY,\hat\beta_1)=\sigma^2 c^T\mathbf{w}$ where $w_i=(x_i-\xbar)/S_x^2$; since
$c^T\mathbf{w}=\frac{1}{S_x^2}(\sum c_ix_i-\xbar\sum c_i)=0$, they are uncorrelated. (More
generally $c^T\bY=c^T\bep$ is a residual direction and $\hat\bb$ lives in the column space.)

\smallskip\noindent\textbf{P5.} Here
$\bX=\begin{pmatrix}1&0&0\\1&1&0\\1&0&1\end{pmatrix}$, which has rank $3$ --- full rank ---
so \emph{everything} is estimable and $\hat\bb=\bX^{-1}\bY$. Solving:
$\hat\alpha_1=Y_1$, $\hat\alpha_2=Y_2-Y_1$, $\hat\alpha_3=Y_3-Y_1$. Hence
\[
  \widehat{\alpha_2-\alpha_3}=Y_2-Y_3,\qquad \Var=2\sigma^2 .
\]
(Sanity check via $a^T=c^T\bX$ with $a=(0,1,-1)$: $c=(0,1,-1)$ works, and indeed
$c^T\bY=Y_2-Y_3$.)

\vfill
\begin{center}\small\itshape
Companions: Simple Linear Regression;
Quadratic Forms, Nested Models and Testing; One-Way ANOVA.
\end{center}

\end{document}
